%!TEX root = ../thesis.tex
% Chapter: Methodology

\chapter{Data, Preprocessing, Methods, and Experimental Protocol} \label{sec:methodology}

The pipeline runs from raw-mammogram preprocessing through patch extraction and labelling to feature computation, classification, and evaluation.

\section{Dataset: INbreast} \label{sec:dataset}

We use INbreast \cite{moreira2012inbreast}, a public full-field digital mammography database of 410 DICOM images (cranio-caudal and medio-lateral-oblique views) from a breast centre in Porto. Three metadata sources are linked by a numeric \texttt{file\_id}: the \texttt{INbreast.xls} sheet (BI-RADS assessment \cite{sickles2013acr}, ACR density 1--4, laterality, view); OsiriX XML annotations (343 of 410 images) giving region-of-interest polygons for lesions; and LabelMe-format pectoral-muscle masks \cite{aliniya2023pectoral} for 202 MLO views.

Labels come directly from this metadata. For the \textbf{tissue} task, ACR 1--2 maps to non-dense (class 0) and ACR 3--4 to dense (class 1); for the \textbf{mass} task, BI-RADS 2--3 is benign and 4--6 malignant, with BI-RADS 1 excluded. This yields 409 usable images for the tissue task: one image (\texttt{file\_id 53582540}) drops out because its ACR field is blank in the spreadsheet, and as it carries no lesion annotation the exclusion is inconsequential. The tissue classes are moderately imbalanced (282 non-dense (ACR 1: 136; ACR 2: 146) versus 127 dense (ACR 3: 99; ACR 4: 28), a $69\%/31\%$ split), so the majority-class baseline is $\sim$69\%, against which the $\sim$84\% WS-only and $\sim$95.9\% scattering-covariance accuracies (Sections~\ref{sec:stage1},~\ref{sec:scov_tissue}) sit 15 and 27~pp above. Metrics are reported on the natural distribution, without class-weight balancing.

\section{Preprocessing} \label{sec:preprocessing}

Each image passes through four stages, shown end-to-end for one example in Figure~\ref{fig:preproc_walkthrough}:
\begin{enumerate}
    \item \textbf{Normalisation.} The pixel array (read with pydicom \cite{pydicom2024}) is clipped at its 1st and 99th percentiles and rescaled to $[0,1]$. This approximates the paper's DICOM$\rightarrow$PNG min-max windowing while being robust to the few extreme pixels per image (calibration markers, sensor hot spots) that a plain min-max stretch lets dominate the rescaling; percentile clipping raised 10-fold CV accuracy on the tissue task from $\sim$80\% to $\sim$83\%. All 410 files are \texttt{MONOCHROME2} with no VOI LUT or window metadata, so no windowing or inversion is applied.
    \item \textbf{Pectoral-muscle removal (MLO/ML views).} Where a ground-truth JSON mask exists, the pectoral region is zeroed so it is not mistaken for dense tissue; the LabelMe set covers 202 of the $\sim$206 oblique views, and the four uncovered views are handled by the breast-mask step below. The pectoral-mask source has only a small effect, audited in Section~\ref{sec:annot_audit}.
    \item \textbf{Breast-mask segmentation.} Otsu thresholding \cite{otsu1979threshold} on the uint8 image, followed by morphological closing/opening (11-pixel elliptical kernel), extraction of the largest connected component and hole filling, isolates the breast from the background.
    \item \textbf{Contrast enhancement.} CLAHE \cite{pizer1987adaptive} (clip limit 2.0, $8 \times 8$ tiles, via OpenCV \cite{bradski2000opencv}) is applied within the breast mask.
\end{enumerate}

\begin{figure}[h]
\centering
\includegraphics[width=\textwidth]{../data/outputs/report_figures/preprocessing_pipeline.png}
\caption[The preprocessing pipeline on an example mammogram]{The preprocessing pipeline applied to one example MLO mammogram. The raw DICOM (0) is percentile-normalised (1); the ground-truth pectoral-muscle mask (2, yellow) is zeroed out (3); Otsu thresholding removes the non-breast background (4) via the breast mask (5); and CLAHE enhances local contrast within the breast (6). All patches are later cropped from panel~6.}
\label{fig:preproc_walkthrough}
\end{figure}

Preprocessed images and patches are stored as NumPy \texttt{.npy} arrays rather than PNG; an 8-bit-PNG round-trip changed no held-out result (feature-vector MAE $\sim$0.02\%), so this is a practical choice, not an accuracy one.

\textbf{Differences from the reference paper.} Table~\ref{tab:preproc_comparison} summarises them. The most consequential is contrast enhancement: the paper uses a density-aware histogram adjustment (their reference [73]) that adapts to the ACR category, whereas our density-agnostic CLAHE forgoes exactly the density-specific contrast being classified, a likely source of the WS-tissue gap. The remaining differences (PNG conversion; seeded region growing versus ground-truth polygon masks) are tabulated.

\begin{table}[h]
\centering
\caption[Comparison of the preprocessing pipeline: paper versus replication]{Preprocessing pipeline comparison between the reference paper and this replication.}
\label{tab:preproc_comparison}
\begin{tabular}{p{3.0cm}p{5.2cm}p{5.2cm}}
\toprule
\textbf{Step} & \textbf{Paper \cite{razali2023cnn}} & \textbf{This replication} \\
\midrule
Normalisation & DICOM $\rightarrow$ PNG (implicit min-max windowing); rescale to $[0,1]$ & 1st/99th-percentile clip; rescale to $[0,1]$ \\
\addlinespace
Image format & DICOM $\rightarrow$ PNG conversion & Operate directly on DICOM arrays \\
\addlinespace
Background removal & Morphological operations; background set to NaN (MATLAB convention) & Otsu threshold + morphology + largest component + hole filling; background set to 0 \\
\addlinespace
Pectoral muscle removal & Seeded region growing (MLO views) & Ground-truth LabelMe JSON polygon masks \cite{aliniya2023pectoral} (202 MLO views) \\
\addlinespace
Contrast enhancement & Density-aware histogram adjustment [73] (adapts to ACR category) & CLAHE (clip limit = 2.0, tile grid $8 \times 8$); density-agnostic \\
\addlinespace
Patch tissue labelling & Otsu threshold on remaining breast region & Otsu threshold + ACR-guided class filtering \\
\bottomrule
\end{tabular}
\end{table}


\section{Patch Extraction and Tissue Labelling} \label{sec:patch_extraction}

Background patches for the tissue task are sampled from the 343 annotated images (Figure~\ref{fig:patch_sampling}):
\begin{enumerate}
    \item \textbf{Mass exclusion.} All \texttt{Mass} ROIs are rasterised and dilated by 24 pixels to keep background patches clear of lesions.
    \item \textbf{Tissue segmentation.} Within the allowed region (breast mask minus pectoral muscle minus dilated mass), an Otsu threshold splits fatty from fibroglandular tissue, with a 0.03 gap around the threshold to reduce boundary ambiguity.
    \item \textbf{ACR-guided sampling.} Unlike the paper, the ACR score restricts which classes an image may contribute (ACR 1 $\rightarrow$ fatty only, ACR 4 $\rightarrow$ fibroglandular only, ACR 2--3 $\rightarrow$ both), preventing the mislabelling that a forced Otsu split causes on near-single-tissue breasts.
    \item \textbf{Random sampling.} Up to three patches per allowed class per image are drawn at random positions and sizes (14--30\% of the shorter breast dimension, capped at 160--512 px), each requiring full containment, $\geq$65\% target-tissue purity, and pixel standard deviation $\geq$0.035; overlaps are forbidden. All draws use a single seeded generator (\texttt{np.random.default\_rng(34)}), so the set is deterministic.
    \item \textbf{Resizing.} Accepted patches are resized to $224 \times 224$ (area interpolation) and saved as float32.
\end{enumerate}

\begin{figure}[h]
\centering
\includegraphics[width=0.6\textwidth]{../data/outputs/report_figures/patch_sampling_example.png}
\caption[Background-patch sampling with mass exclusion and ACR-guided labelling]{Background-tissue patch sampling on one breast. Lesion ROIs are dilated into an excluded margin (orange dashed); within the remaining region, an Otsu split with the ACR-guided rule assigns patches as fatty (cyan) or fibroglandular (lime), drawn at varied positions and sizes with no overlap.}
\label{fig:patch_sampling}
\end{figure}

This yields 434 background patches (251 fatty, 183 fibroglandular) versus the paper's 397; the difference reflects our ACR-guided sampling and the absence of a hard stopping criterion. Table~\ref{tab:patch_counts} compares the counts.

\begin{table}[h]
\centering
\caption[Number of extracted patches: paper versus this replication]{Number of extracted patches. The paper uses $n = 112$ source images from INbreast (its Table~3 total, spanning both tasks); our pipeline processes all 343 images with XML annotations.}
\label{tab:patch_counts}
\begin{tabular}{llcc}
\toprule
\textbf{Task} & \textbf{Class} & \textbf{Paper \cite{razali2023cnn}} & \textbf{This replication} \\
\midrule
\multirow{2}{*}{Background tissue} & Fatty            & 210 & 251 \\
                                   & Fibroglandular   & 187 & 183 \\
\midrule
\multirow{2}{*}{Mass}              & Malignant        & 76  & 75 \\
                                   & Benign           & 36  & 41 \\
\midrule
                                   & \textbf{Total}   & \textbf{509} & \textbf{550} \\
\bottomrule
\end{tabular}

\vspace{0.3em}
{\small Mass patches are used in the benign-vs-malignant task reported in Section~\ref{sec:results}. The BI-RADS mapping used is 2--3 $\rightarrow$ benign and 4--6 $\rightarrow$ malignant; the paper does not state its exact cut-off, which accounts for the 41/36 difference in benign counts.}
\end{table}


\section{Wavelet Scattering Transform} \label{sec:ws}

The wavelet scattering transform gives a non-learned representation that is locally translation-invariant and stable to small deformations \cite{mallat2012invariant, bruna2013invariant}. For an input $x$, the zeroth-order coefficient is the local average $S_0 x = x \star \phi_J$ (low-pass filter $\phi_J$ at scale $2^J$); the first and second orders are
\[
S_1 x = |x \star \psi_{\lambda_1}| \star \phi_J, \qquad
S_2 x = ||x \star \psi_{\lambda_1}| \star \psi_{\lambda_2}| \star \phi_J,
\]
with $\psi_{\lambda}$ oriented wavelets; concatenating the orders gives the feature vector. We compute this with Kymatio \cite{andreux2020kymatio} at $J=6$ (invariance scale $2^J=64$~px), $L=5$ Morlet orientations \cite{morlet1982wave}, maximum order 2, on $224 \times 224$ patches, and spatially average each coefficient map to one scalar per path, a 406-dimensional vector. Table~\ref{tab:ws_config} sets this against the paper's configuration; the count differs (406 versus 391) only because Kymatio uses all six scales with $L=5$ where MATLAB reserves one scale for the low-pass filter and uses $L=6$ (both follow $1 + JL + \binom{J}{2} L^2$ for the respective $J,L$), a library convention, not a modelling choice.

\textbf{Order and invariance scale.} Both follow Bruna and Mallat \cite{bruna2013invariant}. Scattering energy decays exponentially with path length (over 99\% within $m\le3$), and second-order coefficients recover the higher-moment texture that first-order coefficients discard (enough to separate textures with identical power spectra), so order 2 is standard and higher orders add negligible energy. The large invariance scale suits the tissue task, which is stationary-texture discrimination: Bruna and Mallat show the optimal invariance scale then approaches the patch width (minimising estimator variance), so a large $J$, and the near-complete pooling it produces, is appropriate here.

\begin{table}[h]
\centering
\caption[Wavelet scattering configuration: paper versus this replication]{Wavelet scattering configuration: reference paper (MATLAB) vs.\ this replication (Python / Kymatio).}
\label{tab:ws_config}
\begin{tabular}{lcc}
\toprule
\textbf{Parameter} & \textbf{Paper \cite{razali2023cnn}} & \textbf{This replication} \\
\midrule
Library / toolbox         & MATLAB \texttt{waveletScattering} & Kymatio \cite{andreux2020kymatio} \\
Wavelet type              & Morlet                            & Morlet \\
Invariance scale          & 64 ($2^J$ pixels)                 & 64 ($2^J$ pixels) \\
$J$ (max decomposition)   & 6                                 & 6 \\
Effective wavelet scales  & 5 (1 reserved for low-pass)       & 6 (all used) \\
$L$ (orientations)        & 6                                 & 5 \\
Max scattering order      & 2                                 & 2 \\
Input size                & $224 \times 224$                  & $224 \times 224$ \\
Spatial pooling           & Spatial averaging                 & Spatial averaging \\
Output features           & 391                               & 406 \\
Feature standardisation   & None (raw coefficients)           & \texttt{StandardScaler} \\
\bottomrule
\end{tabular}
\end{table}


\section{Classifier: Subspace k-NN Ensemble} \label{sec:classifier}

Following the paper we classify with a Subspace $k$-nearest-neighbours ensemble \cite{ho1998random, cover1967nearest}: 80 learners, each fitting a 1-NN classifier on a random 190-dimensional subspace of the feature space (seed 34), combined by soft vote. The paper searched 40--100 learners and subspace dimensions 170--210 and selected 80 and 190, which we adopt. This low-capacity, distance-based ensemble suits our small samples: Bruna and Mallat note that such classifiers outperform more flexible discriminative ones (e.g.\ kernel SVMs) when training data are scarce \cite{bruna2013invariant}, which is especially relevant for the $n\approx115$ mass set. Features are standardised with \texttt{StandardScaler} \cite{pedregosa2011scikit} before classification (k-NN is scale-sensitive); the paper does not mention standardising the WS features, so this is a deviation that may affect results.

\textbf{Train--test split.} Patches are split 80/20, stratified by class, at the image level (on unique \texttt{file\_id}) rather than the patch level, so all patches from one mammogram stay in the same partition; for the tissue task this gives 344 training and 90 test patches. Cross-validation uses \texttt{StratifiedGroupKFold} \cite{pedregosa2011scikit} with \texttt{file\_id} as the group. The rationale (and the leakage issue that prompted image-level splitting) is given in Chapter~\ref{sec:results}.


\section{Mass ROI Extraction} \label{sec:mass_rois}

For the mass task, ROIs come from the INbreast XML annotations (Section~\ref{sec:dataset}); the paper refers only to ``the mass detection algorithm'' without specifying it, so we use the dataset's expert annotations directly. Each \texttt{Mass} bounding box is cropped from the already-preprocessed image (inheriting the same normalisation, masking, pectoral removal and CLAHE), resized to $224 \times 224$, and saved as float32. Labels follow the BI-RADS assessment (2--3 $\rightarrow$ benign, 4a--6 $\rightarrow$ malignant). This yields 116 mass patches (75 malignant, 41 benign) against the paper's 112 (76, 36); the small discrepancy, our retention of all \texttt{Name="Mass"} ROIs (including multi-ROI images) and the paper's apparent inclusion of some asymmetry/distortion findings, is reconciled in Section~\ref{sec:mass_count_discussion}.


\section{CNN Branch: ResNet18 Fine-Tuning and Feature Extraction} \label{sec:cnn}

The paper's CNN branch (its Model~1) fine-tunes an ImageNet-pretrained ResNet18 \cite{he2016deep} (chosen over deeper networks because a shallow model overfits less on a few hundred patches \cite{yu2019abnormality}) with the final layer replaced by a $512 \to 2$ projection. We replicate this in PyTorch (\texttt{resnet18(weights=IMAGENET1K\_V1)}), fine-tuning all layers (matching the paper's MATLAB default), replicating each single-channel patch to three channels with ImageNet normalisation, and applying no augmentation. Training follows the paper's Table~4 exactly: Adam ($\eta=10^{-3}$, $\beta=(0.9,0.999)$), $L_2$ weight decay $10^{-4}$, cross-entropy loss, 60 epochs, batch size 80. We then extract the 512-dimensional average-pool features via a forward hook, standardise them, and classify with the same Subspace k-NN ensemble (Section~\ref{sec:classifier}).

\textbf{Validation set.} As a diagnostic not in the paper, we carve a $\sim$20\% image-level validation set from the training images to track per-epoch overfitting (Section~\ref{sec:results}); these patches are excluded from gradient updates but re-included for k-NN feature extraction, so the classifier sees the same training set as the paper.


\section{Final evaluation protocol and reporting} \label{sec:eval_protocol}

The development narrative in Section~\ref{sec:results} reports each stage on a single seed-34 image-level split with group-aware 10-fold cross-validation. This is not enough for the headline comparison: the mass held-out set is only ${\sim}25$ patches (${\sim}24$ on average across the five seeds), so one extra misclassification shifts accuracy by ${\sim}4$~pp and a single split can swing $10$~pp between seeds. The final protocol therefore follows four principles:
\begin{enumerate}
    \item \textbf{Repeated, image-level splits.} Each method is evaluated over five stratified splits (seeds $\{1,7,34,42,99\}$), grouped by \texttt{file\_id} so no mammogram's patches cross the train/test boundary; the spread across seeds exposes the split sensitivity a single number hides.
    \item \textbf{Leakage-free feature extraction.} Any trained extractor (the CNN) is retrained inside each split on that split's training images only.
    \item \textbf{Training-only scaling.} \texttt{StandardScaler} is fitted on the training fold alone.
    \item \textbf{AU-ROC-led reporting.} As the mass task is imbalanced ($74$ malignant vs.\ $41$ benign; a trivial ``always malignant'' rule scores the ${\sim}68\%$ majority baseline), we lead with the threshold-independent AU-ROC and report accuracy, precision, recall and F1 alongside it, all as mean~$\pm$~standard deviation over the five seeds.
\end{enumerate}

Because all seeds resample the same ${\sim}115$ mass patches, the standard deviation indicates split sensitivity rather than a tight confidence interval. Cross-validation is not a headline metric: the repeated holdout already gives an out-of-sample estimate, and CV on learned CNN features is optimistic because the extractor has seen the inner-validation images. For reproducibility, the Python, NumPy and PyTorch seeds are pinned together and cuDNN is placed in deterministic mode; seed~34 is the canonical diagnostic seed, and the learned WS-map CNN head (trained on only ${\sim}90$ patches) is additionally corroborated by a group-aware cross-validation rather than a single split (Chapter~\ref{sec:discussion}).


% ============================================================
% 3. RESULTS
% ============================================================
